\paragraph{Б.3. Функции}
\subparagraph{Б.3.1}
Пусть $A$ и $B$ — конечные множества, а $f : A \to B$ — некоторая функция. Покажите, что
\begin{itemize}
    \item[а.] если $f$ является инъекцией, то $|A| \le |B|$;
    \item[б.] если $f$ является сюръекцией, то $|A| \ge |B|$.
\end{itemize}

\textbf{Ответ:}
\begin{itemize}
    \item[а.] По определению, если функция является инъекцией то для любых $x_1, x_2 \in A$
    таких что $f(x_1) = f(x_2)$, следует, что $x_1 = x_2$. Это означает,
    что разные элементы из $A$ отображаются в разные элементы из $B$, и, следовательно, $|A| \le |B|$.
    \item[б.] Если $f$ является сюръекцией, то для любого $y \in B$ существует хотя бы один $x \in A$, такой что $f(x) = y$.
    Это означает, что все элементы из $B$ "покрыты" элементами из $A$, и, следовательно, $|A| \ge |B|$.
\end{itemize}

\subparagraph{Б.3.2}
Пусть имеется функция $f(x) = x + 1$. Будет ли она биекцией, 
если ее область определения и область значений — натуральные числа $\mathbb{N}$? 
А если ее область определения и область значений — целые числа $\mathbb{Z}$?

\textbf{Ответ:}

Для натуральных чисел $\mathbb{N}$: Функция $f(x) = x + 1$ не является биекцией.
Хотя она инъективна, она не сюръективна, так как у наименьшего числа из $\mathbb{N}$ (будь то $0$ или $1$) 
нет прообраза в области определения (им должно было бы быть число $-1$ или $0$ соответственно).

Для целых чисел $\mathbb{Z}$: Функция $f(x) = x + 1$ является биекцией.
Она инъективна ($x_1 + 1 = x_2 + 1 \Rightarrow x_1 = x_2$) и
сюръективна, так как для любого $y \in \mathbb{Z}$ существует прообраз $x = y - 1$, который также принадлежит $\mathbb{Z}$.

\subparagraph{Б.3.3}
Дайте естественное определение обратного к бинарному отношению. 
(Если отношение является биекцией, то определение должно давать обратную функцию.)

\subparagraph{Б.3.4 $\star$}
Приведите пример биекции $\mathbb{Z} \to \mathbb{Z} \times \mathbb{Z}$.